\section{Introduction}
\label{sec:introduction}

Many event streams are \emph{row-complete but relation-incomplete}: the public
release preserves one row per transaction, visit, referral, or service action,
yet suppresses the key that says which rows belong to the same underlying
episode.  Examples include patients moving through a referral episode,
household members observed through separate visits, maintenance records from a
shared service run, and travelers sharing a vehicle trajectory.  This is not
ordinary missing-data uncertainty.  Every released row may be accurate, while
the relational object needed for an aggregate query is absent.

A standard pipeline first reconstructs one linkage or clustering and then treats
it as observed.  That workflow can produce a precise answer whose precision is
entirely an artifact of the chosen reconstruction.  Our goal is different:
\emph{discover only the aggregate conclusions that survive every latent event
structure consistent with the released information}.  Given public rows,
possibly set-valued timestamps, and a declared simultaneous-capacity bound, we
optimize a downstream aggregate over all admissible ordered-event
decompositions.  The resulting frontier separates conclusions forced by the
release from conclusions that depend on an arbitrary point reconstruction.

The temporal structure is richer than pair linkage or clique clustering.  A
valid event can contain an $A$--$B$--$C$ chain: $A$ overlaps $B$, $B$ overlaps
$C$, and $A$ never overlaps $C$.  Event size may therefore exceed instantaneous
capacity.  What matters is positive-overlap connectivity across time together
with an occupancy bound at every instant.  This distinction arises naturally
in vehicle runs, referral chains, rotating service crews, and other episodes
whose membership changes sequentially.

We introduce \methodname, a release-aware partial-identification and
optimization framework for this setting.  It has three layers.  First, a
feasible-world model maps the public release into all ordered event partitions
that satisfy timestamp support, connectivity, capacity, and core--buffer
conditions.  Second, an interval formulation exploits consecutive-ones
structure to solve the fixed-root, fixed-span run-pricing problem exactly by
linear programming.  Third, Dantzig--Wolfe decomposition and branch-and-price
combine those run columns into globally feasible event partitions and certify
integer aggregate endpoints.

Our contributions are as follows.
\begin{enumerate}[leftmargin=*,itemsep=1pt,topsep=2pt]
    \item \textbf{Certified aggregate discovery.}  We define sharp,
    assumption-indexed frontiers for row-complete, relation-incomplete event
    streams.  Expanding timestamp supports or relaxing capacity nests the
    feasible worlds and weakly widens every fixed-estimand frontier.
    \item \textbf{Structure-exploiting exact algorithms.}  For one rooted run,
    interval connectivity and occupancy yield an integral consecutive-ones LP
    pricing oracle.  The cross-run master is not generally integral, so we add
    exact buffer-usage and Ryan--Foster branching while retaining the same
    structural pricing oracle at every node.
    \item \textbf{Truth-based and public-data evidence.}  Across 3,000
    controlled instances, complete candidate universes cover the true
    aggregate in every case.  A feasible temporal point reconstruction makes a
    wrong threshold decision in roughly $17$--$19\%$ of instances; every such
    error is marked ambiguous by the certified frontier.  Public Chicago and
    New York City releases then illustrate three distinct uncertainty sources:
    latent linkage, timestamp support, and simultaneous capacity.
\end{enumerate}

Ride pooling is our audit instance, not the scope of the method.  The public
city data contain no partner or run ground truth, so we do not claim to recover
co-riders, vehicle runs, realized capacity, or a population effect.  The city
experiments instead ask which aggregate statements are invariant over the
explicitly declared candidate universe and release semantics.  Controlled
truth experiments separately evaluate coverage, point-reconstruction error,
and sensitivity to candidate recall.
